% Pursuing Stacks (cotes 134-2 to 134-8): archive cote 134-4, site page 57 (PDF page 58), typescript page 145.
% typescript. Draft compared: ps3.tex:2904-2969.
% AI-transcribed by Claude Opus 5.5 (claude-opus-5-5, adaptive thinking, effort high), Message Batches API, batch ps2 (msgbatch_01NVwr1Yuo8qJf5rykeGBXPK), 2026-11-12; not yet checked by eye.
\scanpage{134-4/57}
\mspage{145}
we get the familiar felationship
\[ \text{(2)} \qquad \overline{\mathrm{Hom}}(X,Y)\hand{\xrightarrow{\sim}\pi_{\mathrm{o}}}(\underline{\mathrm{Hom}}(X,Y)) \ , \]
where $\overline{\mathrm{Hom}}$ denotes homotopy classes of maps,  with respect to the canonical
homotopy structure of M.

The first example which I have in mind is the case M=(Cat).  I was
a little short yesterday about $W_{(\mathrm{Cat})}$ being ``geometric''  -  the condition
a') of page 142 is evident indeed in terms of the \inserted[typed]{(geometric !)} definition of weak equi-%
valence in terms of non commutative cohomology (where we need to care only
about zero dimensional cohomology ! ).  But condition b) is a consequence
of the fact that O-connectedness in (Cat) can be checked using only $\hand{\Delta}_1$
(i.e. $\hand{\Delta}_1$ is a generating contractor for the canonical homotopy structure
of (Cat)),  and that $\hand{\Delta}_1$ is $\hand{W}_{(\mathrm{Cat})}$-aspheric over the final category e ,
i.e. that for any C in (Cat),  the projection
\[ \mathrm{C}\mathrm{x}\hand{\Delta}_1 \hand{\longrightarrow} \mathrm{C} \]
is in $W_{(\mathrm{Cat})}$ .  I want to start being attentive from now on about what
exactly are the formal properties \xout{xx} of $W_{(\mathrm{Cat})}$  we are using  -  it real-%
ly seems they boil down to very \xout{littlex,} few,  which we have kept using
without ever having to refer to the ``meaning'' of weak equivalence in
terms of cohomology (except for proving the formal properties we needed,
in terms of the \over{\illegible{1}}{precise} definition  of weak equivalence we have been start\xout{t}-%
ting with from the very beginning).

In terms of the canonical homotopy structure in (Cat),  admitting
$\hand{\Delta}_1$ as a generating homotopy interval (a contractor,  as a matter of fact)\hand{,}
we now get a notion of two functors from a category X to another Y
being homotopic,  and a corresponding notion of homotopy class\inserted{es} of functors
from X to Y,  which are in one-to-one correspondance with the connected
components of the category $\underline{\mathrm{Hom}}$(X,Y) of all functors from X to Y. \inserted{An} elemen-%
tary homotopy from a functor f to another g  (with respect of course to
the basic generating interval $\hand{\Delta}_1$ ,  which is always understood here) is
nothing else but a morphism from f into g.  Thus homotopy classes of func-%
tors are nothing but equivalence classes for the equivalence relation in
Hom(X,Y)  generated by the relation  ``there exists a morphism from f to
g ''.  To consider at all this equivalence relation,  rather than the usual
one of isomorphy between functors,  is rather far from the spirit in which
categories are generally being used  -   as is also the very notion of
\inserted[typed]{a functor which is a} weak equivalence,  which has been our starting point.  Topological motiva-%
tions alone,  it seems,  would induce anyone to introduce such barbaric
looking notions into category theory !

\margin{57}
